\subsection{Anillos fielmente planos}
\label{subsection:0.6.6-fr}

\begin{env}[6.6.1]
\label{0.6.6.1-fr}
Sea $\varphi:A\to B$ un homomorfismo de anillos que hace de $B$ una
$A$-álgebra. Las cinco propiedades siguientes son equivalentes:
\begin{enumerate}
  \item[$a)$] $B$ es un $A$-módulo fielmente plano (es decir,
    $M_{(B)}$ es un funtor \emph{exacto} y \emph{fiel} en $M$).
  \item[$b)$] El homomorfismo $\varphi$ es inyectivo y el $A$-módulo
    $B/\varphi(A)$ es plano.

\oldpage[0\textsubscript{I}]{59}

  \item[$c)$] El $A$-módulo $B$ es plano (es decir, el funtor
    $M_{(B)}$ es \emph{exacto}), y para todo $A$-módulo $M$, el homomorfismo
    $x\to1\otimes x$ de $M$ en $M_{(B)}$ es inyectivo.
  \item[$d)$] El $A$-módulo $B$ es plano y para todo ideal $\mathfrak a$ de
    $A$, se tiene $\varphi^{-1}(\mathfrak aB)=\mathfrak a$.
  \item[$e)$] El $A$-módulo $B$ es plano y para todo ideal maximal
    $\mathfrak m$ de $A$, existe un ideal maximal $\mathfrak n$ de $B$ tal
    que $\varphi^{-1}(\mathfrak n)=\mathfrak m$.
\end{enumerate}

Cuando se cumplen estas condiciones, se identifica $A$ con un subanillo de $B$.
\end{env}

\begin{env}[6.6.2]
\label{0.6.6.2-fr}
Sean $A$ un anillo \emph{local}, $\mathfrak m$ su ideal maximal, $B$ una
$A$-álgebra tal que $\mathfrak mB\neq B$ (lo que ocurre, por ejemplo, cuando
$B$ es un anillo local y $A\to B$ un homomorfismo \emph{local}). Si $B$ es un
$A$-módulo \emph{plano}, $B$ es un $A$-módulo \emph{fielmente plano}. En efecto, como
$\mathfrak mB\neq B$, hay un ideal maximal $\mathfrak n$ de $B$ que contiene
$\mathfrak mB$; como $\varphi^{-1}(\mathfrak n)\cap A$ contiene
$\mathfrak m$ y no contiene $1$, se tiene
$\varphi^{-1}(\mathfrak n)=\mathfrak m$ y se aplica el criterio $e)$ de (6.6.1).
Bajo las condiciones indicadas, se ve, pues, que si $B$ es
noetheriano, lo mismo ocurre con $A$ (6.5.2).
\end{env}

\begin{env}[6.6.3]
\label{0.6.6.3-fr}
Sea $B$ una $A$-álgebra que es un $A$-módulo fielmente plano. Para todo
$A$-módulo $M$ y todo sub-$A$-módulo $M'$ de $M$, se tiene (identificando $M$
con un sub-$A$-módulo de $M_{(B)}$) $M'=M\cap M'_{(B)}$. Para que $M$ sea un
$A$-módulo plano (resp. fielmente plano), es necesario y suficiente que $M_{(B)}$
sea un $B$-módulo plano (resp. fielmente plano).
\end{env}

\begin{env}[6.6.4]
\label{0.6.6.4-fr}
Sean $B$ una $A$-álgebra, $N$ un $B$-módulo fielmente plano. Para que $B$
sea un $A$-módulo plano (resp. fielmente plano), es necesario y suficiente que $N$
lo sea.

En particular, sea $C$ una $B$-álgebra; si el anillo $C$ es fielmente
plano sobre $B$ y $B$ fielmente plano sobre $A$, entonces $C$ es fielmente plano
sobre $A$; si $C$ es fielmente plano sobre $B$ y sobre $A$, entonces $B$ es
fielmente plano sobre $A$.
\end{env}

